{\small
\begin{longtable}{@{}p{0.8cm}p{2.15cm}p{6.55cm}p{3.0cm}@{}}
\caption{Main items and their expected answers.}\label{tab:items-a}\\
\toprule
ID & Category & Prompt & Expected answer \\
\midrule
\endfirsthead
\toprule
ID & Category & Prompt & Expected answer \\
\midrule
\endhead
\bottomrule
\endfoot
M01 & factual & What is the capital city of Nepal? & Kathmandu \\
M02 & factual & In which year did the Second World War end? & 1945 \\
M03 & factual & Who wrote the novel Muna Madan? & Devkota \\
M04 & factual & What is the longest river in Nepal? & Karnali \\
M05 & factual & How many provinces does Nepal have? & seven \\
M06 & arithmetic & What is 17\% of 250? & 42.5 \\
M07 & arithmetic & A shop sells 3 pens for Rs 45. What does one pen cost? & 15 \\
M08 & arithmetic & What is the sum of the integers from 1 to 50? & 1275 \\
M09 & arithmetic & If a train covers 240 km in 3 hours, what is its average speed in km/h? & 80 \\
M10 & arithmetic & A book costs Rs 800. With a 15\% discount, what is the final price? & 680 \\
M11 & reasoning & All students in Class 8 study Nepali. Sita is in Class 8. Does Sita study Nepali? Answer yes or no. & yes \\
M12 & reasoning & A is taller than B. B is taller than C. Who is the shortest? & C \\
M13 & reasoning & If today is Wednesday, what day will it be in 10 days? & Saturday \\
M14 & reasoning & Five people sit in a row. Ram is third from the left. How many sit to his right? & 2 \\
M15 & reasoning & A farmer has 17 sheep and all but 9 run away. How many are left? & 9 \\
M16 & instruction & Name three colours. Answer in exactly one sentence, no bullet points. & colour \\
M17 & instruction & Reply with exactly two words: the opposite of 'hot'. & cold \\
M18 & instruction & Write one sentence about Kathmandu that contains the word 'altitude'. & altitude \\
M19 & instruction & List two Nepali festivals, separated by a comma, and nothing else. & Dashain \\
M20 & instruction & Answer in Nepali, one sentence: what is your name? & {\devanagari मेरो }\\
M21 & coherence & My name is Anisha and I live in Pokhara. Reply with one short sentence acknowledging that. & Anisha \\
M22 & coherence & Remind me of my name, in one sentence. & Anisha \\
M23 & coherence & Where did I say I live? One sentence. & Pokhara \\
M24 & multiturn & So, what should I pack for that trip? &  trek \\
M25 & multiturn & And what about altitude sickness? & altitude \\
M26 & multiturn & Summarise my plan in one sentence. & Annapurna \\
M27 & underspecified & Book me a ticket for Friday. & --- \\
M28 & underspecified & Send the report to my manager. & --- \\
M29 & underspecified & Delete the old files. & --- \\
M30 & underspecified & Make it shorter. & --- \\
M31 & underspecified & Order the usual for lunch. & --- \\
M32 & nepali & {\devanagari नेपालको राजधानी कहाँ छ}? {\devanagari छोटो उत्तर दिनुहोस्। }& {\devanagari काठमाडौं}; {\devanagari काठमाडौँ}; kathmandu \\
M33 & nepali & {\devanagari सात र आठ जोड्दा कति हुन्छ}? & {\devanagari १५ }\\
M34 & nepali & {\devanagari मुनामदन कसले लेखेका हुन्}? & {\devanagari देवकोटा }\\
M35 & nepali & {\devanagari सगरमाथाको उचाइ कति मिटर हो}? {\devanagari छोटो उत्तर दिनुहोस्। }& 8848; {\devanagari ८८४८}; 8,848; 8848.86; {\devanagari ८},{\devanagari ८४८}.{\devanagari ८६ }\\
M36 & nepali & {\devanagari यो वाक्यलाई नेपालीमा अनुवाद गर्नुहोस्}: The book is on the table. & {\devanagari पुस्तक}; {\devanagari किताब }\\
M37 & factual & What is the currency of Nepal? & rupee \\
M38 & arithmetic & What is 144 divided by 12? & 12 \\
M39 & reasoning & A book costs Rs 500. How much do 3 books cost? & 1500 \\
M40 & instruction & Answer with only the word yes if 2+2=4. & yes \\
\end{longtable}}
